The objective of this thesis was to investigate and analyze the colorectal cancer (CRC) care process within the SMPHS, proposing improvements to enhance patient prognosis. This chapter focuses on presenting the main results of the study, discussing its methodology and value, assessing the accomplishment of its objectives, and providing recommendations for related fields.

\section{Main Findings}

\subsection*{Process Discovery}
The process discovery phase aimed to understand the functioning and performance of the CRC care process, with a focus on the key stages of diagnosis and treatment. Initially, the primary goal seemed to be the opportune delivery of treatment. However, as the study progressed, it became clear that preventive medicine could significantly improve prognosis and reduce costs, despite the need for substantial effort and resources to design a CRC screening strategy. This shift in focus led to an emphasis on improving disease prognosis as the new goal of the CRC care process.

The data analysis revealed that the current system is failing to provide opportune diagnosis and treatment to patients. This is likely to cause discomfort and mistrust among patients using the public healthcare system, potentially leading many of them to abandon treatment.

\subsection*{Process Improvement}

It was determined that the process improvement would be designing a screening strategy and validating it with a microsimulation model. While this approach does not resolve all the identified issues in CRC care, it means an improvement in patient prognosis, which is the primary objective.

The program was inspired by previous CRC screening models and aimes to utilize state-of-the-art knowledge about the disease and the SMPHS population. After completion, the model showed promising results for improving patient prognosis. However, it also indicated that this enhancement comes with increased overall costs. It is now up to the authorities to assess the feasibility and value of implementing the screening strategy.

\section{Discussion}

\subsection*{Challenges of the Study}

The MND team played a crucial role in supporting the data collection process. However, the absence of a standardized procedure for accessing government data for research purposes complicated this stage. Despite clear collaboration intentions from the SMPHS, obtaining the data was not as immediate as hoped.

The primary data source, the SIGGES system, provided reasonably clean data for analysis but lacked the level of quality and detail required for reviewing complex cases. Additionally, the uncertainty surrounding the data made it challenging to distinguish between complex cases and data noise, leading to the exclusion of these cases from the analysis.

Furthermore, in the beginning, there was no established methodology for Process Mining in the process discovery phase. Therefore, the researcher spent a considerable amount of time discovering most preprocessing techniques on its own. Discovering the MEDCP methodology later provided reassurance that the prior work was correct but highlighted the unnecessary time lost.

One of the most challenging decisions in this project was whether to utilize an existing simulation model or to develop one from scratch. Despite efforts to obtain other simulation models' code for reference or to gather insights not available in their published papers, these attempts were unsuccessful. This made the researcher hesitant before deciding to proceed with developing the model independently.

Another ongoing challenge is the lack of specific medical information about CRC in Chile. The characteristics of the disease, such as prevalence and sojourn times for each stage, can vary significantly between regions. For the microsimulation, having more detailed medical data from the simulated region would have greatly enhanced the reliability of the outcomes.

\subsection*{How can we improve in the Future?}

The data collection process did not have a clear mistake by the researcher, aside from the approach of attempting to gather all data from a single source. A missed opportunity was not exploring additional data sources, such as private clinics, which routinely collect data on this disease.

The overlooking of complex cases could have been mitigated in two ways. Firstly, with a larger and higher-quality dataset, the researcher could have identified more complex cases and potentially uncovered explanations for these instances. Secondly, with additional medical expertise, the researcher would have had more support in understanding the complexities and nuances of the more intricate cases.

Understanding how to preprocess the available data was challenging due to the lack of a clear methodology. And the MEDCP methodology provided clear guidance on the necessary preprocessing techniques for this type of project. It is recommended for future projects to allocate dedicated time and resources to finding a suitable methodology for achieving the desired objectives.

In retrospect, the process of obtaining other simulation models should have been initiated earlier. A key mistake was not starting this phase of the thesis sooner and not being decisive enough on when to start the project from scratch. Starting this phase earlier would have increased the likelihood of obtaining access to a functional simulation model or its code. Additionally, being more decisive about starting the simulation model could have allowed for more time to its development and validation.

Regarding the lack of medical information of CRC from Chile, the researcher has exhausted all available public data sources. To address this challenge, the simulation's code repository has been made public, allowing for future improvements and updates as new medical data becomes available.

\section{Value and Replicability of this Study}

The discovery, analysis, and improvement strategy of the CRC care process are highly valuable for various stakeholders.

For the SMPHS and other healthcare institutions managing CRC, this study provides a comprehensive analysis of the CRC care process, offering valuable insights for future improvement decisions. While the focus is on CRC care in this specific region, similar analyses could be conducted for other diseases covered by the Explicit Guarantees in Health (GES) program or in different regions of the country with minor adjustments.

Regarding Process Mining, the use of the MEDCP methodology adds credibility to the results and demonstrates its effectiveness in conducting projects related to diseases like CRC. This study adds to the growing body of work that uses Process Mining to identify common disease pathways and provide key insights into care process performance.

Finally, the simulation model shows promise in improving patient prognosis. While further validation by external validators is necessary to establish its credibility, the model successfully achieves its goal of creating a valid screening simulation model. Additionally, the use of microsimulation for this particular case proves highly useful in assessing the feasibility and cost-effectiveness of implementing this public policy project.

\section{Conclusions}

The project of understanding and analysing the CRC care process through raw data was succesful, as it allowed the researcher to propose valid advice to the SMPHS on how to improve. The results of the process discovery and analysis were supported by real-data analysis and performance indicators that identified genuine issues the service is facing.

The microsimulation model produced in this study met the desired objectives. The data used was obtained from reliable sources, and the outcomes were consistent with various screening strategies worldwide.  Furthermore, the use of sensitivity analyses on key parameters further strengthened the reliability of the findings. The results confirm an improvement in prognosis with increased population participation in the screening program. While the overall costs of implementing these procedures rise, it is the responsibility of the authorities to evaluate the cost-effectiveness of this project and decide whether to proceed with its implementation.

\section{Further Work}

\subsection*{SIGGES Data Collection and Accesibility}

As previously mentioned, the researcher appreciates the government's efforts to track the diagnosis and treatment of each GES disease. However, this study uncovered indications of potential issues with the data and its collection process. While access to the data was eventually granted, the process was clumsy and time-consuming. If the government aims to promote healthcare system improvements, streamlining the process for accessing anonymized datasets would facilitate research endeavors and contribute to a better understanding of healthcare dynamics.

\subsection*{CRC Microsimulation Model}

The researcher concludes that the simulation developed for this study is well-suited to the available data, given the constraints of time and resources inherent in a master's thesis project. The simulation adequately captures key aspects of the CRC care process, offering valuable insights and informing recommendations for further research and policy development.

The model utilized in this study will be accessible in a public repository, enabling anyone interested to test, critique, or enhance it. The researcher hopes for ongoing development and refinement of this simulation model, welcoming contributions from the broader community to further improve its accuracy and effectiveness.

\subsection*{Replication for other diseases}

The main recommendation from this study is to conduct similar analyses for all diseases within the Chilean public healthcare system. By examining successful implementations of new strategies in other countries, many of these approaches can be evaluated for their feasibility and cost-effectiveness.

While the focus of this study is on colorectal cancer, its findings highlight the wider potential for preventive measures across the range of the 94 diseases covered under the GES system. It is hoped that each of these conditions will undergo thorough analysis for the possible integration of preventive medicine, with the ultimate goal of enhancing public health outcomes nationwide.
